%====================================================================
%  CHAPTER 1: INTRODUCTION
%====================================================================
\chapter{Introduction}
\label{chap:intro}

This is the introduction of the dissertation. Replace this placeholder
text with your own content.

This dissertation explores methods in machine learning (ML) with a focus on artificial intelligence (AI), the Internet of Things (IoT), long short-term memory (LSTM), and support vector machine (SVM) architectures. A key challenge is overfitting, addressed via regularisation.

Equations are numbered by chapter:
\begin{equation}
  y = \mathbf{w}^\top \mathbf{x} + b,
  \label{eq:linear}
\end{equation}
where $\mathbf{w}$ is the weight vector and $b$ is the bias.

\begin{equation}
  \mathbf{W} = \arg\min_{\mathbf{W}}
    \|\mathbf{Y} - \mathbf{XW}\|_F^2 + \lambda\|\mathbf{W}\|_F^2.
  \label{eq:ridge}
\end{equation}

%--------------------------------------------------------------------
\section{General Background}
\label{sec:background}

Provide broad background on your research area here. The use of
AI in engineering has grown significantly, with convolutional neural
network (CNN) and LSTM models becoming standard tools.

%--------------------------------------------------------------------
\section{Motivation and Contribution}
\label{sec:motivation}

Describe your motivation and key contributions here.

\begin{itemize}
  \item Contribution 1: Description of the first contribution.
  \item Contribution 2: Description of the second contribution.
  \item Contribution 3: Description of the third contribution.
\end{itemize}

%--------------------------------------------------------------------
\section{Outline}
\label{sec:outline}

Chapter~\ref{chap:litreview} reviews related work.
Chapter~\ref{chap:method} presents the proposed methodology.
Chapter~\ref{chap:results} reports experimental results.
Chapter~\ref{chap:conclusion} concludes the dissertation.
